\chapter{طراحی و پیاده‌سازی چارچوب}
\label{chap:framework}

این فصل هسته‌ی پایان‌نامه است: معماری لایه‌ای چارچوب، قراردادهای پایداری، لایه‌های مکان‌یابی، اشیای صفحه، زبان مخصوص دامنه، مدیریت درایور و پنجره‌ها، جداسازی داده و گزارش‌گیری. انگیزه‌ی هر تصمیم، یکی از دشواری‌های واقعی تست رابط کاربری است که در طول پیاده‌سازی مشاهده و مستند شد.

\section{نگاه کلی به معماری}
چارچوب از پنج لایه تشکیل شده که از بالا به پایین با هم صحبت می‌کنند و فقط لایه‌ی پایین‌تر با درایور سروکار دارد:
\begin{itemize}
	\item \textbf{تست‌ها:} سناریوهای رفتاری؛ تنها از لایه‌ی زبان دامنه استفاده می‌کنند و هیچ فراخوانی درایوری ندارند.
	\item \textbf{لایه‌ی زبان دامنه (\lr{\gls{DSL}}):} پوشش معنادار برنامه («واحد استان را فیلتر کن»، «مجموع برش‌های نمودار دایره‌ای را بخوان»).
	\item \textbf{لایه‌ی اشیای صفحه (\gls{PageObject}):} برای هر صفحه/زبانه یک کلاس؛ تنها نهاد همان صفحه را می‌شناسد.
	\item \textbf{لایه‌ی مکان‌یابی:} مکان‌یابی کنترل‌ها با استراتژی‌های چندگانه، نهان‌گاه دسته‌ها
	\gls{HandleCache}
	و جستجوی محدوده‌ای.
	\item \textbf{لایه‌ی درایور:} کارخانه‌ی درایور
	\gls{DriverFactory}،
	برپایی سرور، اجبار پیش‌زمینه و ورودی صفحه‌کلید سطح سیستم‌عامل.
\end{itemize}

این جداسازی مستقیماً از الگوی شیء صفحه می‌آید
\cite{fowler2002}
اما با یک قاعده‌ی سخت‌گیرانه‌تر تکمیل شده است: دایرکتوری تست‌ها در بازرسی خودکار به فراخوانی‌های درایور حساس است و وجود حتی یک فراخوانی خام، ساخت را می‌شکند. به این ترتیب مرز لایه‌ها قرارداد شفاهی نمی‌ماند.

\section{قراردادهای پایداری در برنامه}
شکننده‌ترین بخش تست‌های رابط کاربری، مکان‌یابی است. چارچوب دو قرارداد در سطح برنامه‌ی نمونه تعریف می‌کند که برنامه‌های هدف باید رعایت کنند:
\begin{enumerate}
	\item \textbf{شناسه پایدار
	\gls{StableIdentifier}:}
	هر کنترل تعاملی، نام دامنه‌ای ثابت خود را هم در
	\lr{objectName}
	(که در
	\lr{\gls{UIA}}
	\cite{uia}
	به
	\lr{AutomationId}
	می‌رسد) و هم در
	\lr{accessibleName}
	(که به
	\lr{Name}
	می‌رسد) قرار می‌دهد. مکان‌یابی ابتدا با نام و در صورت نیاز با شناسه انجام می‌شود.
	\item \textbf{قرارداد خوانش
	\gls{ReadoutContract}:}
	عناصر فقط‌خواندنی (نشانگرهای کلیدی، برچسب‌های شمارنده) شناسه‌ی پایدار خود را در
	\lr{objectName}
	نگه می‌دارند اما
	\lr{accessibleName}
	آن‌ها همیشه متن نمایش‌داده‌شده را منعکس می‌کند؛ زیرا هیچ کلید ویژگی‌ای در
	\lr{WinAppDriver}
	مقدار داخلی این کنترل‌ها را نمی‌خواند. نمودارهای گرافیکی نیز داده‌ی خود را در یک فهرست خوانا آینه می‌کنند تا تست‌ها به تصویر دسترسی نداشته باشند و به داده.
\end{enumerate}

با این دو قرارداد، تستِ نمودار و نشانگر \pdash{} که در درخت دسترس‌پذیری اصلاً دیده نمی‌شوند \pdash{} به تست داده‌ی خوانا تبدیل می‌شود و تست‌ها در برابر تغییرات ظاهری مقاوم می‌مانند.

\section{لایه‌ی مکان‌یابی و نهانگاه دسته‌ها}
هر جستجوی ریشه‌ای در
\lr{\gls{UIA}}
کل درخت را بازمی‌گردد؛ در برنامه‌ای با جدول ده‌هزار ردیفی، هر جستجو تا بیش از ده ثانیه طول می‌کشد. سه سازوکار پایداری و سرعت پیاده شد: نخست، برنامه‌ی هدف بخش‌های کوچک نام‌گذاری‌شده (نوار فیلتر، نوار صفحه‌بندی) را در درخت آشکار می‌کند و لایه‌ی مکان‌یابی جستجو را به همان محدوده‌ی کوچک محدود می‌کند (پردازش اندازه‌گیری‌شده: از ۱۱ ثانیه به ۰٫۰۲ ثانیه)؛ دوم، دسته‌ی \lr{(Handle)} هر کنترل پس از نخستین یافتن در نهانگاه
\gls{HandleCache}
با کلید (عنوان پنجره، شناسه کنترل) نگه داشته می‌شود و کار با دسته‌ی نهان‌شده در حد ۰٫۰۱ ثانیه است؛ سوم، نهانگاه در هر تغییر پنجره باطل می‌شود (بخش اتصال فازی).

\section{کارخانه‌ی درایور و اتصال فازی}
\lr{WinAppDriver}
جلسه را به یک پنجره‌ی سطح‌بالا گره می‌زند؛ برنامه‌ی نمونه در جریان ورود، پنجره‌ی ورود را می‌بندد و پنجره‌ی اصلی می‌سازد و همین تغییر، جلسه‌ی قدیم را بی‌اعتبار می‌کند. کارخانه‌ی درایور فرایند برنامه را مالکیت می‌کند (اجرای برنامه، پایان دادن به آن در پایان جلسه) و جلسه را به عنوان پنجره‌ی سطح‌بالای برنامه می‌سازد؛ لایه‌ی زبان دامنه در هر فاز (ورود، پنجره اصلی) دوباره به پنجره‌ی جاری متصل می‌شود؛ به این سازوکار «اتصال فازی»
\gls{PhaseAttach}
گفته می‌شود. علاوه بر این، دو رفتار خطرناک ویندوز مهار شده است: اول اینکه ورودی صفحه‌کلید به پنجره‌ی پیش‌زمینه می‌رسد و اگر برنامه در پیش‌زمینه نباشد، متن تست در پنجره‌ی دیگری (مثلاً یک پنجره گفتگو) تایپ می‌شود؛ کارخانه پیش از هر تعامل، پیش‌زمینه را به‌صورت قطعی به برنامه منتقل می‌کند. دوم اینکه
\lr{send\_keys}
درایور به متن اضافه می‌کند نه جایگزین؛ اشیای صفحه ورودی‌های متنی را پیش از تایپ خالی می‌کنند.

\section{ورودی صفحه‌کلید و اجرای قطعی}
تعامل با چرخاننده‌ها (کمبوها) از مسیر صفحه‌کلید انجام می‌شود: باز کردن پنجره‌ی بازشو با کلیک، تایپ سطح سیستم‌عامل متن گزینه (به‌کارافتادن جستجوی نوع‌نویسی
\lr{Qt})
و تأیید با
\lr{Enter}.
این مسیر نه به مکان‌یابی عناصر پنجره‌ی بازشو نیاز دارد و نه در برابر هم‌نامی گزینه‌ها با سلول‌های جدول می‌شکند. همین الگو برای فیلترهای تاریخ تعمیم داده شد: تاریخ با ارسال هشت رقم از طریق خود درایور روی عنصر تاریخ تنظیم می‌شود و بخش‌های سال، ماه و روز به‌ترتیب پر می‌شوند. هر یک از این روش‌ها با یک تست کاوشی
\lr{(probe)}
به صورت مستقل کشف و در توضیحات کد مستند شده است.

\section{تست با زبان دامنه در برابر کد خام درایور}
برای نشان‌دادن ارزش لایه‌ی زبان دامنه، همان سناریو را در دو سطح می‌بینیم. کد زیر همان تعامل با فراخوانی‌های خام درایور است؛ یعنی آنچه نویسنده‌ی تست بدون چارچوب می‌نویسد:

\begin{latin}
\begin{lstlisting}[language=Python]
# raw driver access - what a tester writes WITHOUT the framework
driver = attach("Iran Signal Tracker")
section = driver.find_element(
    "xpath", "//*[@AutomationId='dashboardFilterSection']")
combo = section.find_element(
    "xpath", ".//*[@AutomationId='dashboardProvinceCombo']")
combo.click(); sleep(0.5)      # popup takes keyboard focus
keyboard.type_text("Tehran"); keyboard.press_enter()
sleep(0.6)                     # cascade settle
driver.find_element(
    "xpath", "//*[@Name='dashboardSearchButton']").click()
kpi = driver.find_element(
    "xpath", "//Edit[contains(@AutomationId,'kpiCells')]")
assert kpi.get_attribute("Name") == "1023 cells"
\end{lstlisting}
\end{latin}

و این همان تست با زبان دامنه‌ی چارچوب:

\begin{latin}
\begin{lstlisting}[language=Python]
dash = signed.dashboard()
dash.province("Tehran").search()
assert dash.kpi("cells") == "1023 cells"
\end{lstlisting}
\end{latin}

تفاوت فقط کوتاه‌تر شدن کد نیست. در نسخه‌ی خام، نویسنده‌ی تست باید بداند مکان‌یابی با کدام استراتژی کار می‌کند، چرا بین کلیک و تایپ مکث لازم است، چرا ورودی باید از مسیر صفحه‌کلید سیستم‌عامل برود، و ده‌ها رفتار ریز زیرساخت که در جدول
\autoref{tab:quirks}
فهرست شده‌اند؛ هر خطای ریز در این جزئیات، شکست نمایشیِ تست است نه شکست برنامه. در نسخه‌ی زبان دامنه، همه‌ی آن دانش یک بار در لایه‌ی اشیای صفحه زیست می‌کند و نویسنده‌ی تست با واژگان دامنه («استان»، «جستجو»، «خواندن نشانگر») کار می‌کند؛ اگر فردا جستجو از مسیر دیگری انجام شود، فقط یک کلاس تغییر می‌کند و هیچ تستی بازنویسی نمی‌شود. به بیان دیگر، نسخه‌ی خام دانش فنی را در همه‌جا پخش می‌کند و نسخه‌ی زبان دامنه آن را یک‌جا متمرکز می‌کند؛ همین تمرکز است که نگهداشت مجموعه‌ی ۲۵ سناریویی را برای یک نفر ممکن می‌کند.

\section{جداسازی پایگاه‌داده در سطح تست}
تست‌ها باید مستقل و بازتولیدپذیر باشند
\cite{meszaros2007}.
هر تست پیش از اجرا، فایل پایگاه‌داده‌ی خود را از یک «فیکسچر ثابت»
\gls{Fixture}
از پیش متعهدشده بازمی‌گرداند و سپس مهاجرت‌های طرح را روی آن اعمال می‌کند. نگاشت تست به فیکسچر در یک فایل پیکربندی است و بازگردانی بی‌قیدوشرط است (نه «فقط اگر تغییر کرده») تا هیچ تستی به حالت جلسه‌ی قبلی وابسته نماند. خروجی‌های
\lr{\gls{CSV}}
نیز در هر تست پاک می‌شوند تا تست‌های گزارش‌گیری به فایل‌های باقی‌مانده از اجرای دیگر تکیه نکنند.

طراحی فیکسچرها عمداً چندپایگاه‌داده‌ای است: هر فیکسچر یک فایل
\lr{SQLite}
متعهدشده با محتوای قطعی است و فایل پیکربندی
\lr{JSON}
تعیین می‌کند هر تست کدام فیکسچر را دریافت کند؛ در حالت پیش‌فرض همه‌ی تست‌ها فیکسچر کامل با ۱۰٬۰۸۰ نمونه می‌گیرند و هر تست می‌تواند با یک نگاشت صریح فیکسچر خود را بازنویسی کند. چارچوب هم‌اکنون دو فیکسچر دارد: کامل و خالی (برای رفتار برنامه با داده‌ی صفر)، و افزودن فیکسچر تازه — مثلاً داده‌ی یک استان خاص برای تست‌های آینده — تنها یک فایل
\lr{.db}
و یک خط در فایل پیکربندی است؛ هیچ تغییری در کد چارچوب لازم نیست. فایل نگاشت چنین است:

\begin{latin}
\begin{lstlisting}[numbers=none]
{
  "default": "seed_full",
  "tests/test_tc018_empty_fixture.py::test_tc018_empty_fixture":
    "empty"
}
\end{lstlisting}
\end{latin}

کلید نخست، فیکسچر پیش‌فرض همه‌ی تست‌ها را تعیین می‌کند و هر کلید دیگر یک پوشش سرتاسری است: شناسه‌ی کامل تست (مسیر فایل + نام تابع) به فیکسچر موردنظر نگاشت می‌شود. این شکست «پیکربندی به جای کد» باعث می‌شود انتخاب داده‌ی هر تست در یک نگاه دیده و تغییرش بدون دست زدن به آزمون انجام شود. فیکسچرها هرگز در محل اجرا باز نمی‌شوند (فقط بازنویسی بایت‌به‌بایت) تا بایت‌های فایل متعهدشده سالم بمانند و نتیجه‌ی هر اجرا دقیقاً بازتولید شود.

\section{گزارش‌گیری و سوییت‌ها}
تست‌ها در فایل‌های مستقل به نام سناریو (از
\lr{TC001}
تا
\lr{TC025})
نوشته شده‌اند و فایل پیکربندی سوییت‌ها، نام هر سوییت (\lr{smoke}، \lr{auth}، \lr{filters}، \lr{reporting}، \lr{dashboard}، \lr{data}، \lr{unit}) را به فهرست فایل‌های تست آن نگاشت می‌کند. اسکریپت اجرا، پیش از اجرا از سلامت ساختمان برنامه و سرور درایور پیش‌پرواز می‌گیرد، مجموعه را اجرا می‌کند و در پایان گزارش
\lr{Allure}
\cite{allure}
را بسازد و روی یک سرور محلی باز کند؛ زیرا باز کردن گزارش از مسیر فایل، نسخه‌های جدید گزارش را می‌شکند. تصویر لحظه‌ی خطا نیز به صورت خودکار به گزارش هر تست شکست‌خورده پیوست می‌شود. فایل نگاشت سوییت‌ها نمونه‌ای واقعی از همین پیکربندی است:

\begin{latin}
\begin{lstlisting}[numbers=none]
{
  "smoke":  ["test_tc001_launch", "test_tc012_login",
             "test_tc002_province_filter", "test_tc009_export_csv"],
  "auth":   ["test_tc011_register", "test_tc012_login",
             "test_tc013_logout", "test_tc015_users_admin"],
  "filters": ["test_tc002_province_filter",
              "test_tc003_city_filter", "..."],
  "reporting": ["test_tc009_export_csv",
                "test_tc010_chart_visibility", "..."],
  "dashboard": ["test_tc020_dashboard_filters",
                "test_tc021_dashboard_pie_vs_csv",
                "test_tc022_dashboard_tech_bar"],
  "data":  ["test_tc008_refresh", "test_tc016_pagination",
            "test_tc019_demo_failure", "..."],
  "unit":  ["test_fixtures_unit"]
}
\end{lstlisting}
\end{latin}

هر نام سوییت به فهرست فایل‌های تست خود نگاشت می‌شود (پسوند
\lr{.py}
اختیاری است) و اسکریپت اجرا با یک نام، آن سوییت را اجرا می‌کند؛ انتخاب سوییت در هنگام اجرا (مثلاً فقط
\lr{auth}
در هر تغییر پوشش احراز هویت) بدون هیچ تغییری در کد ممکن است. این فایل در عمل نقش فهرست وام‌های کیفیت را دارد: هر سناریو دقیقاً در یک سوییت ثبت است و جابه‌جایی آن بین سوییت‌ها یک خط تغییر است.

\section{سوت‌های مستندشدهی زیرساخت}
در طول پیاده‌سازی، رفتارهای غیرمنتظره‌ای از زیرساخت مشاهده و هر یک با تست کاوشی مستند شد. جدول
\autoref{tab:quirks}
خلاصه‌ی آن‌هاست؛ ارزش این جدول در این است که هر یک به یک سازوکار دائمی در چارچوب تبدیل شده است.

\begin{table}[H]
\centering
\caption{سوت‌های مشاهده‌شدهی زیرساخت و پاسخ چارچوب}
\label{tab:quirks}
\begin{tabularx}{\textwidth}{@{}XX@{}}
\hline
رفتار مشاهده‌شده & پاسخ چارچوب \\
\hline
\lr{send\_keys}
به متن اضافه می‌کند & پاک‌سازی پیش از تایپ در اشیای صفحه \\
جلسه به یک پنجره گره می‌خورد & اتصال فازی + باطل‌سازی نهانگاه \\
کلیک به پنجره‌ی غیرپیش‌زمینه می‌رسد & اجبار قطعی پیش‌زمینه \\
جستجوی ریشه‌ای بسیار کند & بخش‌های نام‌گذاری‌شده + نهانگاه دسته‌ها \\
عناصر جدول مجازی می‌شوند & شمارش از برچسب‌ها نه از ردیف‌ها \\
داده‌ی نمودار در درخت دیده نمی‌شود & قرارداد خوانش و فهرست آینه \\
\hline
\end{tabularx}
\end{table}

نمونه‌ی زیر تستی است که تنها با زبان دامنه نوشته شده؛ خوانایی و نبودِ هر فراخوانی درایوری معیار طراحی این لایه است:

\begin{latin}
\begin{lstlisting}[language=Python]
def test_tc021_dashboard_pie_vs_csv(signed):
    net = signed.network()
    net.province("Tehran")
    result = net.export().should_create_file()
    per_operator = count_operators(result.newest_file())

    dash = signed.dashboard()
    dash.province("Tehran").search()
    slices = dash.pie_slices()
    assert {label: count for label, count, _ in slices} \
        == dict(per_operator)
\end{lstlisting}
\end{latin}